\chapter{Spring Boot: The Machine Under Every Service}
\label{ch:springboot}

Seven of the eight deployables in this project are Spring Boot
applications. Before any YAML file can make sense, you need three ideas:
what a \emph{starter} pulls in, what \emph{auto-configuration} decides on
your behalf, and --- most important for everything in Parts~V and VI ---
exactly \emph{how Spring resolves a configuration property} when the same
key is defined in several places at once.

\section{The problem Spring Boot solves}

Classic Spring asked you to wire everything: a servlet container to deploy
into, XML or Java config declaring every bean, and a long list of
version-compatible dependencies. Spring Boot inverts this.
An application is a plain \code{main()} method:

\begin{lstlisting}[language=Java]
@SpringBootApplication
public class CourseServiceApplication {
  public static void main(String[] args) {
    SpringApplication.run(CourseServiceApplication.class, args);
  }
}
\end{lstlisting}

\code{@SpringBootApplication} is three annotations in a trench coat:
\code{@Configuration} (this class may define beans),
\code{@ComponentScan} (find \code{@Service}, \code{@RestController},
\ldots{} in this package and below), and \code{@EnableAutoConfiguration}
--- the interesting one.

\subsection{Starters}

A starter such as \code{spring-boot-starter-web} is an empty jar whose
only job is a curated dependency list: Spring MVC, Jackson, and an
\emph{embedded} Tomcat. ``Embedded'' matters operationally: the server is
inside the jar, so \code{java -jar course-service.jar} is a complete web
server. That single fact is what makes the container images of
Chapter~13 so simple --- there is no Tomcat to install, only a JVM to
provide.

\subsection{Auto-configuration}

At startup, Spring Boot walks a list of candidate configurations
and applies each one whose conditions hold. The conditions read like:

\begin{itemize}
  \item \emph{is this class on the classpath?} (\code{@ConditionalOnClass}
  --- you added the Postgres driver, so a \code{DataSource} is configured)
  \item \emph{has the user defined their own bean already?}
  (\code{@ConditionalOnMissingBean} --- your explicit bean always wins)
  \item \emph{is a property set?} (\code{@ConditionalOnProperty})
\end{itemize}

Auto-configuration is why \code{course-service} contains no code creating
a connection pool, a Kafka producer, or an S3 client wiring --- adding the
dependency plus a few properties \emph{is} the configuration. When
behavior seems to appear from nowhere, run the app with
\code{--debug} once: the ``condition evaluation report'' lists every
auto-configuration applied and skipped, with reasons.

\section{Externalized configuration: the property resolution order}
\label{sec:property-order}

This is the single most load-bearing mechanism in the whole book.
Every trick the Kubernetes manifests play in Part~V works because of it.

A Spring property such as \code{spring.datasource.url} can be defined in
many \emph{property sources}. Spring merges them all into one
\code{Environment}, and when two sources define the same key, the more
specific source wins. Simplified to the sources this project actually
uses, from weakest to strongest:

\begin{enumerate}
  \item defaults inside \code{application.yml} packaged in the jar
  \item config fetched from config-server (Chapter~2)
  \item OS environment variables
  \item command-line arguments
\end{enumerate}

\begin{decision}[why the cluster overrides with env vars]
The same image must run unchanged on a laptop (Compose) and in k3s. Baking
environment differences into the jar would mean one image per environment
--- exactly what containers are meant to avoid. Environment variables
outrank both the packaged YAML and the config-server, so the Kubernetes
Deployment can repoint a service at \code{postgres-course:5432} without
touching a single file the developers own. One artifact, many
environments, differences injected at the edge.
\end{decision}

\subsection{Relaxed binding}

Environment variable names cannot contain dots. Spring therefore maps
\code{SPRING\_DATASOURCE\_URL} $\rightarrow$ \code{spring.datasource.url}:
uppercase, dots become underscores. Read any Deployment in
\code{deploy/k8s/base/} and you can now decode every \code{env:} entry as
``this overrides that YAML key.'' For example
\code{SPRING\_DATA\_MONGODB\_URI} in \code{dictionary-service.yaml}
overrides \code{spring.data.mongodb.uri} from the config-server.

\subsection{Placeholders with defaults}

The project's config files use one more idiom:

\begin{lstlisting}[language=yaml]
app:
  jwt:
    secret: ${JWT_SECRET:mySecretKeyThatIs...}
\end{lstlisting}

\code{\$\{NAME:default\}} means ``use the environment variable if present,
else the default.'' This is why a developer can run auth-service with zero
setup, while production injects the real secret --- and why in
Chapter~16 the Kubernetes Secret only has to provide \code{JWT\_SECRET} to
take control of this value.

\section{The project's file, annotated}

Every service carries only a tiny \code{application.yml}; the interesting
config lives on the config-server. Here is course-service's, complete:

\begin{lstlisting}[language=yaml]
spring:
  application:
    name: course-service            (*@\co{1}@*)
  config:
    import: optional:configserver:http://localhost:8888   (*@\co{2}@*)
eureka:
  client:
    service-url:
      defaultZone: http://localhost:8761/eureka/          (*@\co{3}@*)
\end{lstlisting}

\begin{description}
  \item[\co{1}] The service's identity. The config-server uses it to pick
  which file to serve (\code{configurations/course-service.yml}); Eureka
  uses it as the registration name; the gateway routes to it as
  \code{lb://course-service}. One string ties all three systems together.
  \item[\co{2}] Bootstrap-time import of remote config.
  \code{optional:} means ``start anyway if the config-server is down'' ---
  kind in dev, and debatable in production (a service that boots with
  half its config may fail in stranger ways than one that refuses to
  boot). The \code{localhost:8888} default is overridden by
  \code{SPRING\_CONFIG\_IMPORT} in Compose and Kubernetes alike.
  \item[\co{3}] Where Eureka lives --- again a localhost default, again
  overridden per environment by \code{EUREKA\_CLIENT\_SERVICEURL\_DEFAULTZONE}.
\end{description}

\section{Profiles}

A Spring \emph{profile} is a named variant of configuration.
\code{spring.profiles.active=native} on the config-server (Chapter~2) is
the project's only profile use so far, but the mechanism generalizes:
\code{application-prod.yml} is loaded on top of \code{application.yml}
when the \code{prod} profile is active, letting one jar carry several
personalities. Kubernetes reduces the need --- the environment injects
differences instead --- which is why this project leans on env vars, not
profile files, for per-environment change.

\section{Actuator}

\code{spring-boot-starter-actuator} adds operational endpoints under
\code{/actuator}: \code{/health} (aggregated readiness of DB, disk,
custom checks), \code{/info}, \code{/metrics}, \code{/prometheus}.
Two facts matter now:

\begin{itemize}
  \item The starter must be on the classpath. Exposure configuration
  without \code{spring-boot-starter-actuator} is inert --- the project
  shipped exactly that mistake in four services and found it through a
  failing Kubernetes probe (Chapter~17).
  \item Endpoints must be \emph{exposed}
  (\code{management.endpoints.web.exposure.include}) --- the project
  exposes \code{health,info,metrics,prometheus}.
  \item Exposure is not \emph{authorization}: if Spring Security protects
  \code{/actuator/**}, the endpoint exists but answers 401. This
  exact interaction decides which Kubernetes probes each service can use
  in Chapter~17 --- dictionary-service permits actuator and probes
  \code{/actuator/health}; auth-service does not, and must be probed via
  \code{/v3/api-docs} instead.
\end{itemize}

\section{Outside the project}

Two contrasting examples of the same property machinery:

\emph{A Heroku-style twelve-factor app} configures everything by
environment variable --- \code{DATABASE\_URL}, \code{PORT} --- and carries
an \code{application.yml} of nothing but \code{\$\{...\}} placeholders.
Spring's resolution order makes that style native.

\emph{A bank running one artifact through five environments}
(dev/test/uat/preprod/prod) typically pairs a config-server backed by a
git repository with profiles per environment:
\code{course-service-uat.yml} lives in an audited repo, and promotion is a
git merge. The project's classpath-native config-server (Chapter~2) is
the single-developer version of the same shape.

\begin{pitfall}[my override is being ignored]
Symptom: you set an environment variable but the service still uses the
old value. Usual causes, in order of frequency: the variable name does
not map to the key you think (check the relaxed-binding spelling ---
\code{SERVICEURL} in the Eureka key is one word, an irregular name that
must be copied exactly); the value is also set by a \emph{stronger}
source (a command-line arg); or the service reads the key at build time
rather than from the Environment. Diagnose with the actuator
\code{/actuator/env} endpoint (when exposed), which lists every property
source and which one won.
\end{pitfall}

\begin{quiz}
\begin{enumerate}
  \item A key is defined in the jar's \code{application.yml}, in the
  config-server's file for that service, and as an env var in the
  Deployment. Which value wins, and why is that ordering the right one
  for containers?
  \item What three annotations does \code{@SpringBootApplication}
  combine, and which of them explains why deleting an unused dependency
  can change runtime behavior?
  \item \code{\$\{GOOGLE\_CLIENT\_ID:placeholder\}} --- what does this
  resolve to when the variable is unset, and what operational property
  does that give auth-service at startup?
  \item Why does exposing an actuator endpoint not guarantee that a
  Kubernetes probe can use it?
\end{enumerate}
\end{quiz}
